% from ega0-20-fr.tex
\oldpage[0\textsubscript{IV-1}]{140}
En effet, dans le diagramme (20.6.9.1), les deux flèches verticales de
droite sont des homomorphismes bijectifs (20.5.12, (ii)).
\end{proof}

\begin{theorem}[20.6.11]
\label{0.20.6.11-fr}
Soient $u:A\to B$, $v:B\to C$ deux homomorphismes d'anneaux, $L$ un
$C$-module. Supposons vérifiée l'une des conditions suivantes :
\begin{enumerate}
  \item[\rm(i)] $L$ est un $C$-module injectif.
  \item[\rm(ii)] $\Upsilon_{C/B/A}$ est facteur direct du $C$-module
$\Omega^1_{B/A}\otimes_B C$ et
$u_{C/B/A}:\Omega^1_{C/A}\to\Omega^1_{C/B}$ est inversible à droite.
\end{enumerate}
Alors l'homomorphisme canonique (20.6.7.1)
\[
  \operatorname{Exalcom}_{B/A}(C,L)\longrightarrow
  \operatorname{Hom}_C(\Upsilon_{C/B/A},L)
\]
est bijectif.

En particulier, si $\Omega^1_{C/B}$ et $\Omega^1_{C/A}$ sont des
$C$-modules projectifs, l'homomorphisme canonique (20.6.7.1) est
bijectif.
\end{theorem}

\begin{proof}
Le fait que chacune des conditions (i), (ii) entraîne que (20.6.7.1) est
bijectif résulte dans les deux cas de la définition de $\alpha_1$. On
notera d'ailleurs que la condition (ii) est \emph{nécessaire et
suffisante} pour que l'homomorphisme (20.6.7.1) soit bijectif pour
\emph{tout} $C$-module $L$ (Bourbaki, \emph{Alg.}, chap. II,
3\textsuperscript{e} éd., \S~2, n\textsuperscript{o}~1, prop.~1). Si on
suppose que $\Omega^1_{C/B}$ et $\Omega^1_{C/A}$ sont des $C$-modules
projectifs, alors, dans la suite exacte (20.6.1.1),
$\operatorname{Ker}(u_{C/B/A})$ est un $C$-module projectif, car la suite exacte
\[
  0\longrightarrow\operatorname{Ker}(u_{C/B/A})\longrightarrow\Omega^1_{C/A}
  \longrightarrow\Omega^1_{C/B}\longrightarrow0
\]
est scindée, $\Omega^1_{C/B}$ étant projectif; comme
$\operatorname{Ker}(u_{C/B/A})=\operatorname{Im}(v_{C/B/A})$, la suite exacte
\[
  0\longrightarrow\Upsilon_{C/B/A}\longrightarrow
  \Omega^1_{B/A}\otimes_B C\longrightarrow
  \operatorname{Im}(v_{C/B/A})\longrightarrow0
\]
est scindée.
\end{proof}

\begin{corollary}[20.6.12]
\label{0.20.6.12-fr}
Supposons que $C$ soit une $A$-algèbre formellement lisse. Il existe alors
un homomorphisme canonique
\begin{equation}
\tag{20.6.12.1}
\label{0.20.6.12.1-fr}
  \operatorname{Exalcom}_B(C,L)\longrightarrow
  \operatorname{Hom}_C(\Upsilon_{C/B/A},L).
\end{equation}
En outre, cet homomorphisme est bijectif si l'une des conditions (i),
(ii) de (20.6.11) est vérifiée.
\end{corollary}

\begin{proof}
En effet, il résulte de l'hypothèse sur $C$ que
$\operatorname{Exalcom}_B(C,L)=\operatorname{Exalcom}_{B/A}(C,L)$, et
l'homomorphisme (20.6.12.1) n'est autre que (20.6.7.1).
\end{proof}

\begin{corollary}[20.6.13]
\label{0.20.6.13-fr}
Si $C$ est une $A$-algèbre formellement lisse et si $\Omega^1_{C/B}$ est
un $C$-module projectif, l'homomorphisme (20.6.12.1) est bijectif.
\end{corollary}

\begin{proof}
En effet, on sait alors que $\Omega^1_{C/A}$ est un $C$-module projectif
(20.4.9), les topologies étant discrètes.
\end{proof}

\begin{env}[20.6.14]
\label{0.20.6.14-fr}
Les notations restant les mêmes, supposons maintenant en outre que $A$,
$B$, $C$ soient des $\Lambda$-algèbres et $u$, $v$ des
$\Lambda$-homomorphismes, ce qui revient à se donner trois homomorphismes
d'anneaux
\[
  \Lambda\xrightarrow{s}A\xrightarrow{u}B\xrightarrow{v}C.
\]

